\fchapter{Ej 2 Pág 39}

\fsection[\textcolor{waveBlue2}{Valoración:}]{\boldit{\textcolor{waveBlue2}{Valoración:}}{
La extensión de vida del producto es una de las grandes preocupaciones en relación con la protección 
del medioambiente. La Comisión Europea, en contribución a los objetivos del Pacto Verde Europeo, ha 
introducido nuevos derechos de los consumidores en materia de reparaciones que promoverán y facilitarán
la reparación y la reutilización. Lee el artículo publicado por el Centro de Documentación Europea de la 
Universidad de Granada y reflexiona sobre cómo podrían las tecnologías habilitadoras digitales contribuir
al cumplimiento de estas medidas en tu sector profesional.
}}
% Accede al artículo a través de este enlace o descargándolo en el apartado de recursos de la web de Editéx: https://bit.ly/48pnrJU>

\begin{solution}

Tras leer el artículo, veo que la nueva normativa europea de derecho a la reparación va mucho más allá 
de un simple cambio legal: 
\vspace{0.5em}

Es una herramienta real para reducir residuos y alargar la vida de los equipos eléctricos.
En el sector de la electricidad y las instalaciones eléctricas, muchos aparatos (cuadros eléctricos, inversores 
solares, cargadores de vehículo eléctrico, sistemas de iluminación LED, etc.) se desechan antes de tiempo solo 
porque una pieza pequeña ya no está disponible o la reparación resulta demasiado cara o compleja. 
\vspace{0.5em}

La directiva obligará a los fabricantes a ofrecer piezas de repuesto, información técnica y 
condiciones razonables durante años, lo que cambiará por completo la forma de trabajar.
\vspace{0.5em}

Las tecnologías habilitadoras digitales serán clave para aprovechar al máximo esta normativa:
\vspace{0.5em}
\begin{itemize}[label=\textbullet]
	\item{\boldit{La impresión 3D} permitirá fabricar bajo demanda carcasas o soportes que ya no se producen.}

	\item{\boldit{La impresión 3D} permitirá fabricar bajo demanda carcasas o soportes que ya no se producen.
	Los pasaportes digitales de producto y blockchain darán acceso inmediato a esquemas, historiales 
		y guías de desmontaje.}
	\item{\boldit{El IoT y la inteligencia artificial} facilitarán el mantenimiento predictivo
		y el diagnóstico rápido.}
	\item{\boldit{La realidad aumentada} ayudará a realizar reparaciones complejas con mayor 
		precisión y seguridad.}




\end{itemize}

En definitiva, pasaremos de sustituir equipos completos a repararlos de manera eficiente y económica.
Esto no solo cumplirá con los objetivos del Pacto Verde Europeo, sino que reducirá significativamente 
los residuos electrónicos, el consumo de materias primas críticas y las emisiones asociadas, al tiempo 
que ofrecerá a los clientes un servicio más sostenible y profesional. 
\end{solution}
